\documentclass{article}

\begin{document}

\section{Silverleaf Phacelia}

Silverleaf phacelia (\textit{Phacelia hastada}) is a native perennial forb common to the American West disturbed within early-seral foothills up past the timberline, and thriving in disturbed coarse-grained substrates. It possesses a stout taproot upon which several prostrate to erect stems grow 10-30 inches long, with fine silvery hairs covering the shoot structure. Silverleaf phacelia is a popular visitor among pollinators including butterflies, moths, and a variety of bee genera that preferentially visit the forb. Although it is not considered good livestock forage, small mammals such as ground squirrels and birds (notably greater sage-grouse, \textit{Centrocercus urophasianus}), eat silverleaf phacelia flowers as sustenance. It blooms multiple white-to-purple flowers in helicoid cymose inflorescence throughout its long blooming period from mid-spring through late summer. Silverleaf phacelia is grown easily from seed and reliably re-seeds itself after high-intensity fires, and in disturbed or even contaminated soil. For this reason it is referred to as a "workhorse" species of roadside soil stabilization and remediation of acidic and heavy-metal contaminated land.

\end{document}

